\chapter{Estado da arte}
\label{chap:estado-da-arte}

Este capítulo recolle o marco teórico e os antecedentes necesarios para contextualizar a solución
proposta neste traballo. En particular, revísanse os conceptos e técnicas máis relevantes
relacionados coa visión por computador, a aprendizaxe profunda aplicada á análise de imaxes e os
enfoques existentes para a avaliación automatizada do dano en organismos.

\section{Visión por computador aplicada á análise de imaxes biolóxicas}
Nos últimos anos, a visión por computador consolidouse como unha ferramenta de grande utilidade na
análise automática de imaxes biolóxicas, ao permitir a extracción de información cuantitativa a partir
de rexistros visuais de maneira máis rápida, reproducible e escalable ca a inspección manual. Este
tipo de técnicas resultan especialmente valiosas en contextos nos que o volume de datos é elevado ou
nos que a avaliación visual require coñecemento experto, como sucede en tarefas de seguimento
morfolóxico, clasificación de mostras ou detección de anomalías.

No ámbito acuático e mariño, esta tendencia tamén se fixo patente. O traballo de Kumar et
al.~\cite{articuloBiomedical} revisa avances recentes na análise de bioimaxes aplicadas á detección de
deformidades esqueléticas en peixes de acuicultura e en especies empregadas en investigación
biomédica, mostrando como a intelixencia artificial e os métodos modernos de análise de imaxe poden
automatizar procesos tradicionalmente manuais. Aínda que o foco dese estudo non está nas algas nin na
avaliación de dano por herbivoría, si evidencia unha tendencia xeral: a incorporación de técnicas de
visión por computador en problemas biolóxicos permite mellorar a consistencia da análise e ampliar a
capacidade de procesamento en escenarios con datos reais e heteroxéneos.

\section{Aprendizaxe profunda para detección e clasificación}
Dentro deste contexto, a aprendizaxe profunda adquiriu un papel central na análise de imaxes debido á
súa capacidade para aprender representacións complexas directamente a partir dos datos. En
particular, as redes neuronais convolucionais constitúen unha das arquitecturas máis empregadas en
tarefas de clasificación visual, xa que permiten extraer patróns espaciais xerárquicos e adaptarse con
boa eficacia a problemas nos que a aparencia, a textura ou a forma resultan determinantes.

No caso da detección de obxectos, a familia de modelos YOLO converteuse nunha das referencias máis
influintes do campo. A revisión de Ramos e Sappa~\cite{articuloYOLO} analiza unha década de evolución
de YOLO, desde as primeiras versións ata arquitecturas máis recentes, e destaca como trazos
característicos a súa eficiencia, a súa escalabilidade modular e a súa capacidade de adaptación a
dominios diversos. Esta combinación de rapidez e precisión explica a súa ampla adopción en aplicacións
de tempo real e en contornas nas que é necesario localizar automaticamente elementos de interese
dentro da imaxe.

Para o presente traballo, este tipo de modelos resulta especialmente relevante porque a avaliación do
dano non se realiza sobre unha rexión perfectamente recortada e estandarizada, senón sobre imaxes nas
que primeiro é preciso identificar a zona correspondente á alga. Neste sentido, a utilización dun
detector baseado en YOLO encaixa de maneira natural como primeira etapa do pipeline, ao permitir
localizar o obxecto de interese antes de proceder á súa análise posterior.

\section{Regresión ordinal e estimación do dano}
En problemas nos que a variable obxectivo representa niveis ordenados, a formulación mediante
regresión ordinal resulta máis axeitada que unha clasificación multiclase convencional, xa que ten en
conta a relación de orde existente entre as etiquetas. No contexto deste traballo, esta característica
é especialmente relevante, dado que o dano observado sobre as algas non se interpreta como un
conxunto arbitrario de clases independentes, senón como unha progresión de intensidade.

Por este motivo, o uso de enfoques específicos de regresión ordinal permite modelar de maneira máis
fiel a estrutura do problema e reducir erros conceptualmente máis graves, como tratar do mesmo xeito
unha confusión entre niveis próximos e outra entre niveis afastados. Entre estes enfoques, CORAL-CNN
constitúe unha alternativa especialmente interesante pola súa capacidade para manter consistencia entre
as predicións binarias que compoñen a saída ordinal.
\subsection{Regresión ordinal (CORAL-CNN)}
No traballo de Cao et al.~\cite{articuloCoral} formúlase a regresión ordinal como un problema no
que as etiquetas non son simplemente categorías independentes, senón valores cunha orde relativa que
debe ser tida en conta polo modelo. Os autores sinalan que os enfoques baseados en descompoñer esta
tarefa en múltiples subtarefas de clasificación binaria poden producir inconsistencias entre
predicións, o que resulta especialmente problemático cando se pretende estimar unha magnitude
ordenada.

Como resposta a esta limitación, proponse \gls{coral}, un \textit{framework} de regresión ordinal
para redes neuronais que garante consistencia no ranking e puntuacións de confianza monotónicas entre
as distintas subtarefas. Ademais, o artigo destaca que se trata dun método independente da
arquitectura empregada, polo que pode adaptarse a modelos convolucionais comúns sen modificar a súa
estrutura base. A avaliación experimental realizada polos autores no contexto da estimación de idade
mostra unha redución do erro respecto de enfoques ordinais previos, o que reforza a súa utilidade en
problemas nos que a orde entre etiquetas é un aspecto esencial.


\section{Traballos relacionados}
Os traballos previos máis próximos ao presente TFG poden agruparse en dúas liñas principais. Por unha
banda, atópanse os estudos que empregan aprendizaxe profunda para resolver problemas de regresión
ordinal en imaxes. Niu et al.~\cite{articuloRegresionOrdinal} propuxeron un enfoque 
baseado nunha CNN con múltiples saídas para transformar a estimación da idade nun conxunto de
subtarefas binarias correlacionadas, demostrando que a extracción de características e o modelado da
orde poden aprenderse de forma conxunta. Sobre esta base, Cao et al.~\cite{articuloCoral}
formularon CORAL como unha mellora orientada a garantir consistencia entre as saídas ordinais,
reforzando así a adecuación deste tipo de métodos para variables organizadas en niveis ordenados.

Por outra banda, existen traballos centrados na automatización da detección e da cuantificación da
severidade en imaxes biolóxicas. O estudo de Sundar et al.~\cite{articuloPlantDisease} presenta un
sistema para a detección de enfermidades en plantas e a estimación da súa severidade mediante unha
combinación de clasificación con redes convolucionais e análise posterior da área afectada. Aínda que
o dominio de aplicación é diferente, este enfoque resulta relevante porque comparte unha motivación moi
próxima á deste TFG: reducir a subxectividade da avaliación visual experta e proporcionar medidas máis
obxectivas e reproducibles do dano observado.

En conxunto, estes antecedentes mostran que xa existen ferramentas eficaces tanto para a localización
de obxectos en imaxes como para a estimación de variables ordenadas ou de severidade en problemas
biolóxicos. Con todo, a aplicación destas técnicas ao caso concreto da avaliación de mordeduras en
algas laminarias segue sendo un problema sen explorar, especialmente cando se traballa con imaxes
non estandarizadas e cunha formulación que require combinar detección, preprocesado e estimación
ordinal nun mesmo fluxo.

\section{Síntese do capítulo}
O estado da arte revisado pon de manifesto que a visión por computador e a aprendizaxe profunda
ofrecen unha base sólida para abordar problemas de análise automatizada de imaxes biolóxicas. Os
modelos de detección, e en particular a familia YOLO, proporcionan mecanismos eficaces para localizar
obxectos de interese en contornas complexas, mentres que os enfoques de regresión ordinal permiten
modelar adecuadamente variables que representan niveis de severidade ou intensidade.

Tamén se observa que boa parte dos traballos previos se centran en dominios diferentes, como a
estimación da idade facial, a análise biomédica de peixes ou a detección de enfermidades en plantas.
Isto implica que, malia existir antecedentes metodolóxicos relevantes, segue habendo un espazo claro
para adaptar e integrar estas técnicas no contexto específico da avaliación de danos en algas
laminarias. A proposta desenvolvida neste traballo sitúase precisamente nese punto, ao combinar
detección, preprocesado e estimación ordinal para construír un sistema orientado ao apoio da análise
experta sobre datos reais.
